\section{为什么 Generative Modeling 能成为 Unsupervised Learning}

监督学习依赖 $(x,y)$ 标签，而互联网大量数据只有 $x$。生成建模通过预测数据自身构造训练信号，因而可以在没有人工任务标签时学习 representation 和 behavior。关键问题不是“标签为零”，而是自监督 objective 是否迫使模型捕获对下游有用的结构。

\subsection{标签瓶颈与任务扩张}

\term{unsupervised learning}（无监督学习）在课堂中泛指不依赖人工任务标签的表示学习；现代语境常更精确称 \term{self-supervised learning}，因为 target 从原数据自动构造。语言模型的 next-token label 就来自文本本身。

\lecturefigure{slide-10-unsupervised-learning.jpg}{监督任务的数据需求推动模型寻找可从原始数据自动产生的训练信号。}{官方视频 09:09。}

\begin{knowledgebox}{读图：两条路径的区别是标签生产方式}
Supervised learning 需要 task-specific labels，成本高且覆盖窄；unsupervised/self-supervised path 从大规模 raw data 学习通用结构，再通过 prompting 或 fine-tuning 迁移。它并没有消除 human choices：数据收集、过滤、tokenizer、objective 与 evaluation 仍是人为设计。
\end{knowledgebox}

\lecturefigure{slide-11-unlabeled-data.jpg}{互联网提供海量无标签数据，但规模不等于质量或授权。}{官方视频 10:38。}

\begin{knowledgebox}{读图：地球图标代表规模，也代表治理}
大规模 web data 覆盖语言、代码、图像和文化差异，为 pretraining 提供广泛 signal；同时包含 duplication、bias、private data、copyright 与 contamination。课程关注“是否存在数据宝库”，工程上还必须回答来源、许可、清洗、去重、保留与删除。
\end{knowledgebox}

\subsection{Analysis by synthesis 与 autoregressive objective}

\term{analysis by synthesis} 的直觉是：若模型能生成与数据分布相符的样本，它必须学到部分潜在结构。Autoregressive objective 最大化
\[
\mathcal{L}(\theta)=\sum_{t=1}^{T}\log p_\theta(x_t\mid x_{<t}).
\]
每个位置都产生训练信号，适合 scale；但 objective 只要求预测 token，不保证内部 representation 与人类概念一一对应。

\lecturefigure{slide-12-autoregressive-analysis-by-synthesis.jpg}{Autoregressive generation 以预测数据来学习结构，是 analysis by synthesis 的一种实现。}{官方视频 11:11。}

\begin{knowledgebox}{读图：生成是训练接口，representation 是副产物}
模型为了预测下一个 token，会利用语法、主题、实体、情感和世界规律；这些因素可能在 hidden states 中形成可读方向。但模型也可利用 surface shortcut。判断 representation 是否有用，需要 probing、intervention 与 transfer，而不只是低 loss。
\end{knowledgebox}

\subsection{Sentiment neuron：涌现表示与解释边界}

Radford 等人在 character-level language model 中发现某个单元与 review sentiment 高度相关。这个结果说明生成 objective 可能学习可用于分类的 feature，也提醒我们“一个 neuron”通常是特定模型与分析方法下的现象，不应外推为所有概念都单神经元编码。

\lecturefigure{slide-13-sentiment-neuron.jpg}{无监督生成模型中出现与情感相关的表示方向。}{官方视频 13:00；Radford et al. 2017。}

\begin{knowledgebox}{读图：左看分布，右看激活轨迹}
左侧比较正负 review 在 neuron activation 上的分布差异；右侧展示文本位置与 activation 变化。相关性支持该单元携带 sentiment signal，但不能单独证明它是唯一因果机制。可进一步通过 clamping、ablation 或 classifier transfer 检验因果与泛化。
\end{knowledgebox}

\teachervoice{讲者用 sentiment neuron 回答“为什么生成文本会帮助下游任务”：为了 synthesis，模型可能先做 analysis。课堂提醒是可能，不是保证；representation quality 仍需被下游证据验证。}

\subsection{本章小结}

Generative pretraining 用数据本身构造密集监督，能够从无标签规模中学习可迁移表示。它的优势来自信号覆盖，风险来自 objective 与真实任务之间仍有距离。

